\documentclass[10pt, letterpaper]{article}
\usepackage[utf8]{inputenc}
\usepackage{amsmath, amssymb, amsthm}
\usepackage{geometry}
\usepackage{enumitem}
\usepackage{verbatim}
\usepackage{graphicx}
\usepackage{booktabs}

% Adjust margins
\geometry{margin=1in}

% Title
\title{\textbf{Validation of Simplified Logic for \textbf{Claude's Cycles}}}
\author{}
\date{}

\begin{document}

\maketitle

This note presents empirical evidence validating the simplified C implementation of the algorithm for \textit{Claude's Cycles}. While theoretical concerns suggested potential bottlenecks for indices $c=1$ and $c=2$, comprehensive testing confirms that the logic successfully generates Hamiltonian cycles for a generic range of odd numbers $m = 2n + 1$, specifically for $n \in \{1, 2, 3, 4, 5\}$ ($m=3, 5, 7, 9, 11$).

\section{Methodology}

The algorithm operates on a 3D grid of size $m \times m \times m$. The state transition logic is determined by the sum of coordinates modulo $m$ ($s = (i+j+k) \pmod m$):
\begin{itemize}
    \item If $s=0$ or $s=m-1$, the move is constrained to axes $j$ or $k$.
    \item Otherwise, the move is constrained to axis $i$.
\end{itemize}
Three distinct cycle types ($c=0, 1, 2$) are generated by selecting the primary increment direction based on index $0$, $1$, or $2$ of the move string $d$.

\section{Results}

The code was executed generically for $m = 2n + 1$ where $n = 1, 2, 3, 4, 5$. In all cases, the cycle visited exactly $m^3$ unique vertices before returning to the start, confirming a Hamiltonian Cycle.

\subsection{Code Implementation:} \input{note-on-claude-cycles-code}

\subsection{Execution Output:} \input{note-on-claude-cycles-output}

\end{document}

